% !PRINTED 453
\weberpage{453}

Sind die $a$ nicht unabhängige Veränderliche, sondern selbst wieder Functionen von anderen Veränderlichen, so bleibt alles dies gültig, so lange nur die Aenderungen der $a$ nach der Stetigkeit erfolgen.

Wir können also hiernach in allen diesen Fällen von den $m$ Wurzeln einer Gleichung $m^{\text{ten}}$ Grades sprechen, die als Zahlen oder als algebraische Functionen ihrer Co\"efficienten aufzufassen sind.

Wenn nun $\Omega$ ein beliebiger Körper und $f(x)$ eine Function in $\Omega$ ist, so sind zwei Fälle möglich.

Die Function $f(x)$ ist entweder zerlegbar oder nicht zerlegbar in Factoren niedrigeren Grades, deren Co\"efficienten dem Körper $\Omega$ angehören. Im ersten Falle heisst die Function $f(x)$ \gesperrt{reducibel}, im zweiten \gesperrt{irreducibel} in $\Omega$. Man spricht bisweilen auch, wenn der Körper $\Omega$ als bekannt und feststehend betrachtet werden kann, wenn es z.\,B. der Körper der rationalen Zahlen ist, schlechtweg von reducibeln oder irreducibeln Functionen, ohne die Bezeichnung \textquotedbl{}im Körper $\Omega$\textquotedbl{} ausdrücklich hinzuzufügen. Dass aber an sich die genaue Präcisirung des Körpers, auf den sich die Reducibilität bezieht, oder der als Rationalitätsbereich betrachtet wird, sehr wesentlich ist, ergiebt sich aus folgenden Bemerkungen.

Eine lineare Function ist selbstverständlich in jedem Körper irreducibel. Multiplicirt man mehrere Functionen der Form (1) mit einander, so entsteht eine Function derselben Form, die aber dann natürlich reducibel ist, da sie wieder in die Factoren zerlegt werden kann, aus denen sie entstanden ist.

Eine Function $f(x)$, die in $\Omega$ irreducibel ist, kann in einem erweiterten Körper $\Omega'$, der aus $\Omega$ durch irgend eine Adjunction entsteht, reducibel werden; so wird jede Function $f(x)$ reducibel in dem Körper $\Omega'$, der aus $\Omega$ durch Adjunction einer Wurzel $\alpha$ von $f(x) = 0$ entsteht; denn es kann von $f(x)$ ein linearer Factor $x - \alpha$ abgesondert werden. In dem Körper; der aus \gesperrt{allen} Zahlen besteht, ist jede Function mit Zahlenco\"efficienten reducibel. % [sic] semicolon as printed

I. \gesperrt{Eine irreducible Function $f(x)$ kann mit einer anderen Function $F(x)$, deren Co\"efficienten demselben Körper angehören, keinen gemeinsamen Theiler haben, wenn nicht $F(x)$ durch $f(x)$ theilbar ist.}
